Um die Beziehung $\epsilon^{\prime \alpha \beta} = -\epsilon^{\alpha \beta}$ korrekt herzuleiten, müssen wir die Rolle der Determinante der Transformation $A$ berücksichtigen. Die Transformation eines antisymmetrischen Tensors wie des Levi-Civita-Symbols unter einer linearen Transformation $A$ ist gegeben durch:

\begin{align*}
\epsilon^{\prime \alpha \beta} &= \det(A) A^\alpha_{\ \gamma} A^\beta_{\ \delta} \epsilon^{\gamma \delta}.
\end{align*}

Da $\epsilon^{\alpha \beta}$ antisymmetrisch ist, gilt $\epsilon^{\alpha \beta} = -\epsilon^{\beta \alpha}$. Für die inverse Transformation $A^{-1}$, die bei orthogonalen Transformationen $A^{-1} = A^T$ ist, ergibt sich:

\begin{align*}
\epsilon^{\alpha \beta} &= \det(A^{-1}) (A^{-1})^\alpha_{\ \gamma} (A^{-1})^\beta_{\ \delta} \epsilon^{\prime \gamma \delta}.
\end{align*}

Da $\det(A^{-1}) = \frac{1}{\det(A)}$ und für orthogonale Transformationen $\det(A) = \pm 1$, folgt:

\begin{align*}
\epsilon^{\prime \alpha \beta} &= \det(A) A^\alpha_{\ \gamma} A^\beta_{\ \delta} \epsilon^{\gamma \delta} \\
&= \pm A^\alpha_{\ \gamma} A^\beta_{\ \delta} \epsilon^{\gamma \delta}.
\end{align*}

Für eine Drehung in zwei Dimensionen ist $\det(A) = 1$, und für eine Spiegelung ist $\det(A) = -1$. Daher ergibt sich für eine Spiegelung:

\begin{align*}
\epsilon^{\prime \alpha \beta} &= -\epsilon^{\alpha \beta}.
\end{align*}

Dies zeigt, dass die transformierte Version des Levi-Civita-Symbols das negative des ursprünglichen Symbols ist, wenn die Transformation eine Spiegelung beinhaltet.

%CHECK: Berücksichtigung der Determinante der Transformation, explizite Herleitung der Vorzeichenänderung bei Spiegelung.